% Source: round12/stability/project_input.tex
\section{Instantiating the stability transfer with the project's multiplier theorem}
\label{s28:sec:stability-project-input}

The uniform Hadamard-gap input is already available in the current
project, not an additional prime-distribution assumption. Precisely,
\path{../../../papers/hadamard_multibase_final.tex}, the uniform-gap
theorem, proves that for each large integer $x$ there is
$d\le\exp(D\log x/\log\log x)$ such that every integer $y\le x$ has a
Hadamard order in $(y,y+d]$. Its formal counterpart is
\path{../../../lean/HadamardMultibase/CompleteResults.lean},
\path{Complete.quantitative_gap_bound}. We use the theorem as stated;
no files in that proof or its verification manifests are changed.

Here is an even more direct dyadic construction using just the
unconditional multiplier input and the uniform buffer lemma of the
arithmetic branch. For $h=2^s$ large choose
\[
 S(h)=\max\{T(h),\lceil\log_2s\rceil\},\qquad
 g=2^{S(h)},\qquad k=h-g.
\]
Since $T(h)=O(\log h/\log\log h)$, eventually
$S(h)\le s-2$, while $g\ge s$ and
\[
                 g=\exp(O(\log h/\log\log h))=h^{o(1)}.
\]
The integer $q=2^{s-S(h)}-1$ is positive odd and at most $h$.
Its admissible exponent is at most $T(h)\le S(h)$, so
\[
                    k=2^{S(h)}q
\]
is a known Hadamard order. This choice is uniform in all sufficiently
large dyadic $h$ and does not use GRH, prime gaps, or an
irrational-rotation subsequence.

\begin{corollary}[Unconditional near-square instability from the multiplier bound]
\label{s28:cor:stability-project-rate}
At orders $n=kh$ obtained above, there are exactly isotropic uniform
laws on $m=h^2$ atoms, on $n-1$ bits, with
\[
 \begin{split}
 C(\mu)/(1/n)&=1-n^{-1/2+o(1)},\\
 1/n-C(\mu)&=n^{-3/2+o(1)},\\
 m&=n+n^{1/2+o(1)},\\
 \operatorname{dist}_{\rm TV}(\mu,\mathcal H_n)
       &=1-(1+o(1))n^{-1/2}.
 \end{split}
\]
They are strict local collision maxima. Their common additive
R\'enyi entropy excess is $n^{-1/2+o(1)}$.
Both $H_k$ and $H_n$ exist unconditionally from the project's
proved inputs.
\end{corollary}
\begin{proof}
The construction gives $s\le g<h/2$. Choose $g$ nonzero omitted
slopes containing an additive basis and apply the spanning-slope
theorem of the preceding pass. It yields exact isotropy, strict
local maximality, collision $1/h^2$, and exact distance $1-1/k$.
Tensoring $H_k$ and Walsh $H_h$ gives $H_n$.
Since $n\sim h^2$ and $g=h^{o(1)}$, the identities
\[
 1-\frac{C}{1/n}=\frac gh,\quad
 \frac1n-C=\frac{g}{kh^2},\quad
 m-n=hg,\quad
 H_\alpha(\mu)-\log n=\log(h/k)\sim g/h
\]
give all the assertions.
\end{proof}

\paragraph{Formalization.}
The independent \texttt{ProjectChallenge} verifies the full unconditional
corollary for one actual family, with a law at every sufficiently large
dyadic exponent. The arithmetic multiplier theorem is proved in the
included, permissively licensed source; no prime-distribution statement
is assumed. For positive quantities $f_i$ at orders $n_i\to\infty$, the
power notation $f_i=n_i^{a+o(1)}$ is recorded by the exact limit
$\log f_i/\log n_i\to a$. The normalized distance remainder satisfies
$\sqrt{n_i}/k_i\to1$. The same laws satisfy all collision, support,
positive-order R\'enyi, min-entropy, and sharp attained mixture-distance
conclusions. Every definition in the independent specification is concrete.

Thus the project's own unconditional existence theorem gives a
stronger asymptotic stability counterexample than the independent
ordinary-prime-gap route, even assuming RH in that prime route.
The RH rate remains useful as an independently sourced comparison.
The positive all-order existence question is unchanged: these
counterexamples occur at orders where matrices are already known.
They concern full $n-1$-bit isotropy, not the punctured extraction input.

